\section{Introduction}
\label{sec:introduction}

Mobile robots deployed in industrial interiors increasingly need more than a
map: they need environment \emph{knowledge}---which objects exist, where they
are, what they afford---that can be queried in task language, trusted enough
to act on, and kept current as the site changes. Terrestrial laser scanning
(TLS) is an attractive source for such knowledge: survey-grade scanners
deliver millimeter-accurate, globally registered geometry of an entire site
in minutes, and are already routine in construction and facility management.
Yet the path from a raw scan to knowledge a robot can safely use is broken in
three places.

First, \emph{semantic labels do not survive contact with industrial scenes}.
Open-vocabulary 3D classifiers trained on web-scale or indoor-furniture
distributions degrade sharply on industrial content: in our experiments a
state-of-the-art point--text embedding model misclassifies most equipment,
and---counterintuitively---restricting its vocabulary to the classes actually
present makes accuracy \emph{worse} (66.7\% $\rightarrow$ 33.3\%,
Sec.~\ref{sec:exp-ablation}), indicating an embedding-level domain gap rather than a class-list
problem. Pipelines that propagate such labels downstream hand the robot
confident fictions.

Second, \emph{ungrounded language interfaces hallucinate}. Letting an LLM
answer questions about a scene from its own priors, or from an unvalidated
label dump, produces fluent but unsupported claims---20.8\% of atomic factual
claims in our LLM-only benchmark condition (Sec.~\ref{sec:exp-query}). For a robot, a
hallucinated fire extinguisher is not a harmless chat artifact.

Third, \emph{scene models are built once and abandoned}. Most semantic
mapping pipelines produce a single-epoch artifact with no principled update
path: a re-scan either rebuilds everything---losing object identities that
downstream planners and logs refer to---or is not integrated at all. Real
sites drift continuously; in our own test room, roughly one third of object
geometry changed within two weeks (Sec.~\ref{sec:exp-update}).

This paper presents \textbf{TLS-SMF} (TLS-driven Semantic Modeling
Framework), which addresses all three failures with one design principle: \emph{separate what can be made deterministic from what
must remain probabilistic, and control the probabilistic part with explicit
confidence machinery}. Concretely, we contribute:

\begin{enumerate}
\item \textbf{A deterministic geometric and explicit-modeling backbone.}
From registered TLS input, seeded-RANSAC structure removal, fixed-parameter
clustering with a conservative giant-split and a structure-remnant filter,
and PCA-based explicit modeling run fully automatically; same-host re-runs
are byte-identical. Determinism is not an implementation nicety: it is what
makes verification results carryable, diffs meaningful, and update costs
proportional to actual change.
\item \textbf{Confidence-aware multi-view VLM verification with automatic
escalation.} Each object is classified per-view under a forced answer
schema; votes are fused confidence-weighted; split votes escalate
automatically to a wider view set; and objects that still fail to converge
are retained as explicitly \texttt{unverified}---never silently dropped and
never resolved by a human in the loop.
\item \textbf{Identity-preserving incremental knowledge-graph update.}
Node IDs are minted once and carried by two-pass matching (spatial, then
type-constrained move detection); re-scans apply node-level
\texttt{unchanged/updated/moved/inserted/absent} upserts with per-revision
history, and unchanged objects re-verify for free by content-hash carry-over.
\item \textbf{Corrective decomposition and condition-varied detection.}
Diffuseness and wall-compound gates turn ``too messy to be an object''
from a discard decision into an automatic re-splitting trigger (wall-band
stripping, support-plane and panel-bank cuts), and a structure-condition
ensemble re-detects under varied wall/ceiling removal, recovering
wall-flush objects any single condition amputates while exposing label
instability as a per-object attribute.
\item \textbf{Owner-in-the-loop revisions over a single semantic store.}
The room occupant's refutations, restorations, and corrections of labels,
geometry, and attributes enter as highest-trust but time-variant evidence
through the same revision mechanism as re-scans; the management console,
the mission mediator, and an Isaac~Sim twin are all derived from the one
revisioned graph, so a correction propagates to every consumer from a
single write.
\end{enumerate}

We validate the framework on two industrial indoor scenes with
owner-confirmed ground truth---plus a cross-floor cafeteria stress scene
that bounds domain transfer (Sec.~\ref{sec:exp-ablation})---through five
benchmark families: segmentation
repeatability and sensitivity, verification accuracy and escalation cost, a
four-condition knowledge-grounded query benchmark with hallucination traps,
synthetic update scenarios, and a real three-epoch study of one room spanning
a manual scan, an autonomous exploration scan, and a controlled-change
re-scan at 99.6\% coverage. Beyond the headline gains---escalation
+11.1/+6.3 points, gated querying +31 points over LLM-only with 3 of 4
hallucination traps blocked,
re-verification at 3.1\% of rebuild cost---the study surfaces
failure modes we believe are as valuable as the successes: phantom objects
from ceiling soffits, unanimous-but-wrong verifications, clutter absorption
of small low-reflectance objects, and a coverage paradox in which a
\emph{better} scan over-merges density-based clustering. Each is analyzed
and either fixed by a method component or reported as a bounded limitation.

The framework extends our ICTC 2026 multi-view VLM classification study with
escalation, structural gating, identity matching, incremental update, and
the multi-epoch evaluation; the knowledge-graph lineage follows our ISIS 2025
work. Both are cited and the deltas are stated explicitly in
Sec.~\ref{sec:related}.
